\documentclass[12pt,a4paper]{article}
\usepackage[top=21mm,bottom=23mm,left=26mm,right=26mm]{geometry}
\usepackage{amsmath,amssymb,xcolor}
\definecolor{link}{HTML}{205973}
\usepackage[hypertexnames=false,colorlinks=true,linkcolor=link,urlcolor=link]{hyperref}
\hypersetup{pdftitle={Appendix B.3. Sampling from an approximate occupancy law}}
\setlength{\parindent}{0pt}
\setlength{\parskip}{5pt}
\raggedbottom
\newcommand{\R}{\mathbb R}
\newcommand{\Prob}{\mathbb P}
\newcommand{\E}{\mathbb E}
\newcommand{\cF}{\mathcal F}
\newcommand{\TV}{d_{\rm TV}}
\newcommand{\ep}{\epsilon}
\newcommand{\al}{\alpha}
\newcommand{\app}{\pi_{\rm approx}}
\newcommand{\pistar}{\pi_{\mu^*,\nu}}
\newcommand{\lead}[1]{\par\smallskip\noindent\textbf{#1}\enspace}
\begin{document}
{\LARGE\bfseries Appendix B.3. Sampling from an\\[2pt] approximate occupancy law\par}
\medskip
When program (18) draws states from an approximation to the optimal policy's
occupancy, its fitted value function still admits a finite-sample guarantee.
The error can be bounded using either the population total variation or a
population likelihood ratio and the observed optimal slacks. The purely
empirical likelihood-ratio inequality on p.~33 can fail with probability at
least one half, regardless of the coupling.

\section*{Which distribution enters the fit?}
Consider a finite discounted cost-minimizing MDP with discount
$0<\al<1$, costs $g_{x,a}$ and transition probabilities $p(x,y,a)$. Let
$T_\mu J=g_\mu+\al P_\mu J$ be the policy Bellman operator and
$TJ=\min_\mu T_\mu J$ its componentwise minimum. An optimal policy
$\mu^*$ has value $J^*=TJ^*=T_{\mu^*}J^*$. Given an initial distribution
$\nu$, abbreviate the ideal and actual sampling distributions by
\[
 \pi=\pistar=(1-\al)\sum_{j\ge0}\al^j(P_{\mu^*}^{\top})^j\nu,
 \qquad \rho=\app .
\]
Write $q^\top f=\sum_x q_x f(x)$ and
$\|f\|_{1,\nu}=\sum_x\nu_x|f(x)|$. The candidate class $\cF$ consists of
$J(x)=\langle\Phi(x),\mathbf z\rangle+b$, with $\Phi$ taking values in a
real Hilbert space, $\|\mathbf z\|\le C$ and $|b|\le B$.
Put $r_J=(J-T_{\mu^*}J)_+$ and $s_J=(J-TJ)_+$, so $0\le r_J\le s_J$.

Draw $X_1,\ldots,X_N$ independently from $\rho$ and write
$\widehat\rho^\top f=N^{-1}\sum_i f(X_i)$. Program (18) fits $\widehat J$
by maximizing over $J\in\cF$ and $N$ slacks
\[
 \nu^\top J-cN^{-1}\sum_i s_i,\qquad
 s_i\ge s_J(X_i),\qquad c=\frac{2}{1-\al}.
\]
There is an optimizer on every sample. Since $c>0$, every optimal slack
is minimal at its draw. To evaluate the slack away from the draws, use the
canonical statewise extension $s_{\widehat J}$. Denote the observed mean
slack by $\overline s=N^{-1}\sum_i\widehat s_i$.

\lead{The TV coefficient.}
With $\TV(p,q)=\frac12\sum_x|p_x-q_x|$, any $0\le f\le L$ satisfies
\begin{equation}
 (p-q)^\top f=(p-q)^\top(f-L/2)\le L\TV(p,q).
 \label{eq:tv}
\end{equation}
Opposite point masses and their separating indicator attain equality.
Applied to empirical laws, the coefficient in B.3 is therefore
$2/(1-\al)$ for standard TV, or $1/(1-\al)$ for the full $\ell^1$ distance.

\newpage
\section*{A guarantee for every fitted optimizer}
Let $K=\max_x\|\Phi(x)\|$, $D=4CK(1+\al)+4B(1-\al)+2\|g\|_\infty$,
and $L=D/2$. For every $J\in\cF$, $0\le r_J,s_J\le L$.
Indeed $J-T_\mu J$ has feature term bounded by $CK(1+\al)$, offset
term bounded by $B(1-\al)$, and cost term bounded by $\|g\|_\infty$.
A greedy action realizes $TJ$ in each state. The stated $L$ covers all
these terms. Choose $\ep>0$, $0<\delta<1$ and a positive integer $N$ with
\[
 N\ge D^2\bigl(1+\sqrt{2\log(1/\delta)}\bigr)^2/\ep^2 .
\]
Define $d=\TV(\pi,\rho)$. The following bounds refer to the product law
$\rho^{\otimes N}$ of the states actually used in the fit.

\lead{Uniform approximation.}
With probability at least $1-\delta$, every optimizer satisfies
\begin{equation}
 \|J^*-\widehat J\|_{1,\nu}\le
 \frac{3+\al}{1-\al}\min_{V\in\cF}\|J^*-V\|_\infty
       +c(\ep+Ld).
 \label{eq:best}
\end{equation}
In particular the change in this guarantee tends to zero with the
population distance $d$, with $C,B,K$ and the MDP fixed.

\lead{Population comparison.}
Fix a comparator $V\in\cF$ before drawing the states and write
$A_\rho(V)=\nu^\top(J^*-V)+c\rho^\top s_V$. With probability at least
$1-\delta-\delta^4$, every optimizer satisfies
\begin{equation}
 \|J^*-\widehat J\|_{1,\nu}\le A_\rho(V)+c(2\ep+Ld).
 \label{eq:fixed}
\end{equation}
If $\pi_x\le(1+k)\rho_x$ for all $x$ and $k\ge0$, the same event gives
\begin{equation}
 \|J^*-\widehat J\|_{1,\nu}\le
 A_\rho(V)+c\bigl((2+k)\ep+k\overline s\bigr).
 \label{eq:ratio}
\end{equation}
The observable correction $k\overline s$ comes with an additional
$k\ep$ from transporting the fitted function's concentration error.
The domination hypothesis applies without dividing by a zero probability.

\lead{The deterministic comparison.}
For any fit $J$, let $u=J-J^*$. The policy fixed-point equation yields
$u\le r_J+\al P_{\mu^*}u$, hence
$u_+\le r_J+\al P_{\mu^*}u_+$. Summing the discounted iterates gives
$\nu^\top u_+\le\pi^\top r_J/(1-\al)$. Using
$\|J^*-J\|_{1,\nu}=\nu^\top(J^*-J)+2\nu^\top u_+$ and empirical
optimality with the feasible comparator slacks $s_V(X_i)$, we obtain
\begin{equation}
 \|J^*-\widehat J\|_{1,\nu}\le
 \nu^\top(J^*-V)+c\bigl(\widehat\rho^\top s_V+\eta\bigr)
 \quad\text{if }\pi^\top r_{\widehat J}\le\widehat\rho^\top r_{\widehat J}+\eta .
 \label{eq:path}
\end{equation}
Indeed, optimality bounds $-\nu^\top\widehat J+c\widehat\rho^\top
s_{\widehat J}$ by the same expression at $V$, and $r_{\widehat J}\le
s_{\widehat J}$. This step applies to the optimizer of the displayed sampled program.

\newpage
\section*{Proof under the actual product law}
Theorem~3, applied to the fixed policy $\mu^*$ under $\rho$, gives a single
event of probability at least $1-\delta$ on which
\begin{equation}
 \rho^\top r_J\le\widehat\rho^\top r_J+\ep
 \qquad\text{for every }J\in\cF .
 \label{eq:uniform}
\end{equation}
The uniform quantifier allows substitution of the observed fit
$\widehat J$. Applying \eqref{eq:tv} before \eqref{eq:uniform} gives
$\eta=\ep+Ld$ in \eqref{eq:path}. Under the ratio assumption,
nonnegativity of $r_{\widehat J}$ gives the alternative estimate
\begin{align*}
 \pi^\top r_{\widehat J}
 &\le(1+k)\rho^\top r_{\widehat J}
 \le(1+k)(\widehat\rho^\top r_{\widehat J}+\ep)\\
 &\le\widehat\rho^\top r_{\widehat J}
         +k\overline s+(1+k)\ep .
\end{align*}
Thus we may also take $\eta=k\overline s+(1+k)\ep$.

For \eqref{eq:best}, put $e=\|J^*-V\|_\infty$. The $\al$-contraction of
$T$ gives $s_V(x)\le(1+\al)e$ pointwise, and
$\nu^\top(J^*-V)\le e$. Substitution of these bounds into \eqref{eq:path}
gives coefficient $1+c(1+\al)=(3+\al)/(1-\al)$ for $e$.
This holds for all $V$ on the same event. The value vectors in $\cF$ form
a compact set. Projecting the coefficient vector onto the span of the
finitely many $\Phi(x)$ preserves every evaluation, and the radius
constraints in this span and on the offset are compact. The minimum
in \eqref{eq:best} is attained.

For \eqref{eq:fixed} and \eqref{eq:ratio}, use $0\le s_V\le D/2$ for the
one comparator $V$ fixed in advance. Changing one draw changes
$\widehat\rho^\top s_V$ by at most $D/(2N)$. If $D>0$, the
bounded-differences inequality in particular gives
\[
 \Prob\!\left(\widehat\rho^\top s_V-\rho^\top s_V\ge\ep\right)
 \le \exp(-2N\ep^2/D^2)\le\delta^4 .
\]
The last inequality uses
$N\ep^2\ge D^2(1+\sqrt{2\log(1/\delta)})^2
\ge 2D^2\log(1/\delta)$. If $D=0$, the residual vanishes identically.
Intersecting the comparator event with \eqref{eq:uniform} costs at most
$\delta+\delta^4$ in failure probability. Inserting each choice of
$\eta$ into \eqref{eq:path} proves the two bounds. A comparator that
minimizes $A_\rho$ can be chosen before drawing the sample because the
admissible value vectors are compact and $A_\rho$ is continuous.

\lead{What closeness of the laws controls.}
Empirical TV also depends on how the samples are paired. Two independent
single draws from the same fair two-state law have empirical TV one with
probability $1/2$. For fixed $N$, one can instead pair each coordinate by
putting $\min(\pi_x,\rho_x)$ on the diagonal at state $x$. Pair the
residual masses by their normalized product when their common total
$d>0$. This gives mismatch probability $d$ per draw. Independent copies
give identical samples with probability $(1-d)^N$. The population-TV term
in \eqref{eq:best}--\eqref{eq:fixed} requires only the actual sample.

\newpage
\section*{Why the empirical ratio can fail}
Use states $0,1$, one available action, absorbing transitions and
$\al=1/2$. Take $g_1=-1$, $g_0=1$, $\Phi=0$, and strictly positive radii
$C=B=1$. Every admissible fit is the constant $J=b$, $|b|\le1$.
Its canonical statewise slack is
\[
 s_J(1)=(1+b/2)_+=1+b/2,\qquad
 s_J(0)=(-1+b/2)_+=0 .
\]
Consequently $s_J=a\mathbf 1_{\{1\}}$ with $a\ge1/2$, including for
every data-dependent optimizer. The factor $a$ cancels from the claimed
empirical inequality. All discounted occupancies in this model equal their
initial distributions, so we may realize any desired ideal law.

Choose $\pi_1=1-t$, $\pi_0=t$, with $0\le t\le1$, and $\rho_1=\rho_0=1/2$.
Then $\pi_x\le 2\rho_x$ for both states, so the appendix's ratio
parameter can be taken to be $k=1$. Choose an odd $N$, and let
$Y\sim\pi^{\otimes N}$, $X\sim\rho^{\otimes N}$ on any common
probability space. Write $A$ for the event that every ideal draw is $1$,
and $B$ for the event that fewer than half of the approximate draws are
$1$. The first marginal gives $\Prob(A)=(1-t)^N$. Complementing every
bit in a fair sample exchanges $B$ and $B^c$ because $N$ is odd, giving
$\Prob(B)=1/2$. Every coupling therefore satisfies
\[
 \Prob(A\cap B)\ \ge\ (1-t)^N-\tfrac12
 \ \ge\ \tfrac12-Nt .
\]
On $A\cap B$ the ideal mean of $s_{\widehat J}$ equals $a$, and its
approximate mean is less than $a/2$. Thus
\[
 (\widehat\pi-\widehat\rho)^\top s_{\widehat J}
       > k\widehat\rho^\top s_{\widehat J}.
\]
For $t=0$, failure is exactly $B$, of probability $1/2$ at every odd
sample size, with the same two population laws throughout. If both
distributions must be strictly positive, set $t=1/(4N)$. Failure
probability is then at least $1/4$ at every odd $N$. In this latter
family the ideal distribution varies with $N$. Neither obstruction
depends on choosing an unfavorable coupling or an unfavorable optimizer.

{\small\raggedright
\lead{Sources.} Bhat, Farias, Moallemi and Zheng,
\href{https://moallemi.com/ciamac/papers/kernel-adp-2012.pdf\#page=32}{Appendix~B.3, pp.~32--33}.
Formal proofs cover
\href{https://github.com/nikhilpb/kadp/blob/main/lean/SamplingLaw.lean}{the transport bounds and empirical ratio counterexample},
\href{https://github.com/nikhilpb/kadp/blob/main/lean/WeightedSampling.lean}{the joint infimum over fits and weights}, and
\href{https://github.com/nikhilpb/kadp/blob/main/lean/SamplingCoupling.lean}{the empirical coupling and equal-law example}.\par}
\end{document}
